% Part: second-order-logic
% Chapter: metatheory
% Section: second-order-arithmetic

\documentclass[../../../include/open-logic-section]{subfiles}

\begin{document}

\olfileid{sol}{met}{spa}

\olsection{د دويمې درجې حساب}

په ياد ولرئ چې د پېانو د حساب تيوري~$\Th{PA}$ د~$\Th{Q}$ اته
بديهي اصول رانغاړي:
\begin{align*}
& \lforall[x][\eq/[x'][\Obj 0]]\\
& \lforall[x][\lforall[y][(\eq[x'][y'] \lif \eq[x][y])]]\\
& \lforall[x][(\eq[x][\Obj 0] \lor \lexists[y][\eq[x][y']])]\\ 
& \lforall[x][\eq[(x + \Obj 0)][x]]\\
& \lforall[x][\lforall[y][\eq[(x + y')][(x + y)']]]\\
& \lforall[x][\eq[(x \times \Obj 0)][\Obj 0]]\\
& \lforall[x][\lforall[y][\eq[(x \times y')][((x \times y) + x)]]]\\
& \lforall[x][\lforall[y][(x < y \liff \lexists[z][\eq[(z' + x)][y]])]]\\
\intertext{او ورسره د لاندې بڼې ټولې جملې:}
& (!A(\Obj 0) \land \lforall[x][(!A(x) \lif !A(x'))]) \lif \lforall[x][!A(x)].
\end{align*}
وروستۍ بڼه يو «سکيما» دے: داسې قالب چې د حساب د ژبې نامتناهي ډېرې
!!{sentence}s جوړوي، د هر !!{formula}~$!A(x)$ دپاره يوه۔ دې قالب ته
د استقرا (د لومړۍ درجې) \emph{بديهي سکيما} وايو۔ د پېانو په
\emph{دويمې درجې} حساب~$\Th{PA^2}$ کښې استقرا په يوه جمله بيانېدے
شي۔ $\Th{PA^2}$ د پورته لومړنيو اتو بديهي اصولو او د (دويمې درجې)
\emph{استقرا له بديهي اصل} څخه جوړه ده:
\[
\lforall[X][((X(\Obj 0) \land \lforall[x][(X(x) \lif X(x'))]) \lif \lforall[x][X(x)])].
\]
دا وايي: که د !!{domain} يو فرعي سټ~$X$، $\Assign{\Obj{0}}{M}$ ولري
او له هر~$x \in \Domain{M}$ سره $\Assign{\prime}{M}(x)$ هم ولري
(يعنې «د جانشين تر عمل لاندې تړلے» وي)، نو د !!{domain} ټول څيزونه
لري (يعنې $X = \Domain{M}$)۔
\textbf{د ژباړې د نسخې سمون (OLSOL-001):} په سرچينه کښې د دې
مساوات د رياضيکي نښو دننه يوه زياته تړونکې قوسه راغلې وه؛ دلته
قوسه د توضيحي عبارت برخه ده۔

د استقرا بديهي اصل تضمينوي چې هر !!{structure} چې دا اصل پوره کوي،
په~$\Domain{M}$ کښې يوازې هغه !!{element}s لري چې بديهي اصول ئې
شتون غواړي؛ يعنې د~$\num{n}$ ارزښتونه د $n \in \Nat$ دپاره۔
د~$\Th{PA^2}$ هېڅ مدل نامعياري شمېرې نۀ لري۔

\begin{thm}
\ollabel{thm:sol-pa-standard}
که $\Sat{M}{\Th{PA^2}}$، نو $\Domain{M} =
\Setabs{\Value{\num{n}}{M}}{n \in \Nat}$.
\end{thm}

\begin{proof}
دې $N = \Setabs{\Value{\num{n}}{M}}{n \in \Nat}$ وي، او فرض کړئ
$\Sat{M}{\Th{PA^2}}$۔ څرګنده ده چې د هر $n \in \Nat$ دپاره
$\Value{\num{n}}{M} \in \Domain{M}$؛ نو $N \subseteq \Domain{M}$۔

اوس د بل لوري شمول وښيو۔ داسې د متغير ګومارنه~$s$ وګورئ چې
$s(X) = N$۔ د فرض له مخې:
\begin{align*}
\Sat{M}{\lforall[X][((X(\Obj 0) \land \lforall[x][(X(x)
      \lif X(x'))]) \lif \lforall[x][X(x)])]}, & \text{ نو}\\
\Sat{M}{(X(\Obj 0) \land \lforall[x][(X(x) \lif X(x'))]) \lif
  \lforall[x][X(x)]}[s]. & 
\end{align*}
د دې شرطي جملې مقدم وګورئ۔ $\Value{\Obj{0}}{M} \in N$، نو
$\Sat{M}{X(\Obj{0})}[s]$۔ دويم عطف~$\lforall[x][(X(x) \lif X(x'))]$
هم صدق کوي۔ ځکه فرض کړئ $x \in N$۔ د~$N$ له تعريفه، د کوم~$n$
دپاره $x = \Value{\num{n}}{M}$۔ له دې څخه
$\Assign{\prime}{M}(x) = \Value{\num{n+1}}{M} \in N$۔ نو
$\Assign{\prime}{M}(x) \in N$۔

نو $\Sat{M}{X(\Obj 0) \land \lforall[x][(X(x) \lif X(x'))]}[s]$ لرو۔
له دې امله $\Sat{M}{\lforall[x][X(x)]}[s]$۔ دا وايي چې د هر
$x \in \Domain{M}$ دپاره $x \in s(X) = N$؛ نو $\Domain{M}
\subseteq N$۔
\end{proof}

\begin{cor}
\ollabel{cor:sol-pa-categorical}
د~$\Th{PA^2}$ هر دوه مدلونه آيزومورف دي۔
\end{cor}

\begin{proof}
د \olref{thm:sol-pa-standard} له مخې، د~$\Th{PA^2}$ د هر مدل تعريف ساحه
د~$\Value{\num{n}}{M}$ له ارزښتونو پرته بل څه نۀ لري۔ هر داسې مدل
د~$\Th{Q}$ مدل هم دے۔ د \olref[mod][mar][stm]{prop:thq-standard} له
مخې، هر داسې مدل معياري دے؛ يعنې له~$\Struct{N}$ سره آيزومورف دے۔
\end{proof}

پورته مو $\Th{PA^2}$ هغه تيوري وبلله چې لومړني اته حسابي بديهي
اصول او د دويمې درجې د استقرا بديهي اصل لري۔ په حقيقت کښې د دويمې
درجې منطق د بيان د ځواک له برکته، د دويمې درجې پېانو حساب دپاره
يوازې \emph{لومړني دوه} حسابي بديهي اصول او استقرا بس دي۔

\begin{prop}\ollabel{prop:sol-pa-definable}
دې $\Th{PA^{2\dagger}}$ هغه د دويمې درجې تيوري وي چې لومړني دوه
حسابي بديهي اصول (د جانشين بديهي اصول) او د دويمې درجې د استقرا
بديهي اصل لري۔ نو $\le$، $+$ او $\times$ په~$\Th{PA^{2\dagger}}$
کښې تعريفېدے شي۔
\end{prop}

\begin{proof}
د دې ښودلو دپاره چې $\le$ تعريفېدے شي، داسې فورمول~$!A_\le(x, y)$
موندل غواړو چې $\Sat{N}{!A_\le(\num n, \num m)}$ هله او يوازې هله
چې $n \le m$۔ دا فورمول وګورئ:
\[
!B(x, Y) \ident Y(x) \land \lforall[y][(Y(y) \lif Y(y'))]
\]
څرګنده ده چې د~$Y \subseteq \Nat$ سټ له خوا $!B(\num n, Y)$ هله او
يوازې هله پوره کېږي چې $\Setabs{m}{n \le m} \subseteq Y$؛ نو
$!A_\le(x, y) \ident \lforall[Y][(!B(x, Y) \lif Y(y))]$ اخيستلے شو۔

د دې ليدلو دپاره چې جمع تعريفېدے شي، پام وکړئ چې $k+l=m$ هله او
يوازې هله چې داسې تابع~$u$ شته چې $u(0)=k$، د هر $n$ دپاره
$u(n')=u(n)'$، او $m = u(l)$۔ دا معادلت د لاندې فورمول په وسيله
په~$\Th{PA^{2\dagger}}$ کښې د جمع د تعريف دپاره کارولے شو:
\[
!A_+(x,y,z) \ident \lexists[u][(u(\Obj 0)=x \land \lforall[w][u(w')=u
 (w)'] \land u(y)=z)]
\] 
\textbf{د ژباړې د نسخې سمون (OLSOL-002):} د جانشين د شرط دننه
د کلي کميت ايښودل شوي متغير لپاره سرچينه د ناسمې نښې کاروي؛
دلته هماغه کميت ايښودل شوے متغير په دواړو ځايونو کښې راغلے دے۔
څرګنده بايد وي چې $\Sat{N}{!A_+(\num k, \num l, \num m)}$ هله او
يوازې هله چې $k+l=m$۔
\end{proof}

\begin{prob}
د \olref[sol][met][spa]{prop:sol-pa-definable} ثبوت بشپړ کړئ۔
\end{prob}

\end{document}
